% Part: computability
% Chapter: recursive-functions
% Section: primes

\documentclass[../../../include/open-logic-section]{subfiles}

\begin{document}

\olfileid{cmp}{rec}{pri}
\olsection{ప్రధాన సంఖ్యలు}

సహజమైన ప్రమేయాలు, సంబంధాలు ఆదిమ పునరావృత్తాలని చూపడానికి పరిమిత
పరిమాణీకరణ, పరిమిత కనిష్ఠీకరణ మనకు చాలా సాధనాలను ఇస్తాయి. ఉదాహరణకు,
``$x$, $y$ను భాగిస్తుంది'' అనే సంబంధాన్ని పరిశీలిద్దాం; దాన్ని
$x\mid y$ అని రాస్తాం. $y$ను $x$తో భాగించినప్పుడు శేషం రాకపోతే,
అంటే $y$ అనేది $x$ యొక్క పూర్ణాంక గుణితమైతే, $x\mid y$ నెరవేరుతుంది.
(ఇది నెరవేరకపోతే, అంటే $y$ను $x$తో భాగించినప్పుడు శేషం $>0$ అయితే,
$x\nmid y$ అని రాస్తాం.) మరో విధంగా, ఏదో~$z$కు $x\cdot z=y$
అయినప్పుడే $x\mid y$. అటువంటి $z$ ఉంటే అది తప్పకుండా $\leq y$.
కాబట్టి ఏదో $z\le y$కు $x\cdot z=y$ అయినప్పుడే $x\mid y$.
$=$, గుణకారం నుంచి పరిమిత అస్తిత్వ పరిమాణీకరణతో ఈ $x\mid y$
సంబంధాన్ని ఇలా
నిర్వచించవచ్చు:
\[
x \mid y \defiff \bexists{z \leq y}{(x \cdot z) = y}.
\]
అందువల్ల $x\mid y$ ఆదిమ పునరావృత్త సంబంధమని చూపాం.
\sourcecorrection{OLTECRREM-003}{ఆంగ్ల మూలంలోని కుండలీకృత వివరణ $x$ను $y$తో భాగించిన శేషాన్ని పేర్కొంది; $x\mid y$ నిర్వచనానికి అనుగుణంగా $y$ను $x$తో భాగించిన శేషంగా సరిచేశాం.}

సహజ సంఖ్య~$x$, $0$ గానీ $1$ గానీ కాకుండా, $1$తోనూ తనతోనూ మాత్రమే
భాగించబడితే దాన్ని \emph{ప్రధాన సంఖ్య} అంటాం. మరో విధంగా, $y\mid x$
అయినప్పుడల్లా $y=1$ లేదా $y=x$ అయ్యే సంఖ్యలే ప్రధాన సంఖ్యలు. $x$
ప్రధానమో కాదో పరీక్షించడానికి $y\le x$ అయిన ప్రతిసంఖ్యకూ $y\mid x$
అవుతుందో లేదో మాత్రమే చూడాలి; ఎందుకంటే $y>x$ అయితే స్వయంచాలకంగా
$y\nmid x$. అందువల్ల, $x$ ప్రధానమైనప్పుడే నెరవేరే సంబంధం
$\fn{Prime}(x)$ను ఇలా నిర్వచించవచ్చు:
\[
\fn{Prime}(x) \defiff x \geq 2 \land \bforall{y \leq x}{(y \mid x \lif y
  = 1 \lor y = x)}
\]
కనుక అది ఆదిమ పునరావృత్త సంబంధం.

ప్రధాన సంఖ్యలు $2$, $3$, $5$, $7$, $11$, మొదలైనవి. ఈ క్రమంలో
$x$వ ప్రధాన సంఖ్యను తిరిగి ఇచ్చే ప్రమేయం $p(x)$ను పరిశీలిద్దాం;
అంటే $p(0)=2$, $p(1)=3$, $p(2)=5$, మొదలైనవి. (సౌలభ్యం కోసం
$p(x)$ను తరచూ $p_x$ అని రాస్తాం: $p_0=2$, $p_1=3$, మొదలైనవి.)

$x$కన్నా పెద్దదైన మొదటి ప్రధాన సంఖ్యను తిరిగి ఇచ్చే ప్రమేయం
$\fn{nextPrime}(x)$ మనకు ఉంటే, $p$ను ఆదిమ పునరావృత్తితో సులభంగా
నిర్వచించవచ్చు:
\begin{align*}
  p(0) & = 2\\
  p(x+1) & = \fn{nextPrime}(p(x))
\end{align*}
$\fn{nextPrime}(x)$ అనేది $y>x$ అయి $y$ ప్రధానమయ్యే అతి చిన్న $y$
కనుక, దాన్ని అపరిమిత అన్వేషణతో సులభంగా గణించవచ్చు. కానీ యూక్లిడ్
ఫలితం వల్ల పరిమిత కనిష్ఠీకరణతో కూడా నిర్వచించవచ్చు: $x$, $\fact{x}+1$
మధ్య ఎప్పుడూ ఒక ప్రధాన సంఖ్య ఉంటుంది.
\[
  \fn{nextPrime}(x) =
  \bmin{y \leq \fact{x}+1}{(y > x \land \fn{Prime}(y))}.
\]
కాబట్టి $\fn{nextPrime}(x)$, తద్వారా $p(x)$ కూడా కేవలం గణనీయాలే
కాక ఆదిమ పునరావృత్త ప్రమేయాలని తెలుస్తుంది.
\sourcecorrection{OLTECRREM-004}{ఆంగ్ల మూలం రెండు చోట్ల ఆర్గుమెంటును ప్రమేయపు పేరు లోపల ఉంచి $\fn{nextPrime(x)}$ అని రాసింది; అధ్యాయం అంతటా వాడిన ప్రమేయ-ప్రయోగ సంకేతనానికి అనుగుణంగా $\fn{nextPrime}(x)$గా సరిచేశాం.}

(ఆసక్తి ఉంటే, యూక్లిడ్ ఫలితానికి సంక్షిప్త నిరూపణ ఇది. $x=0$ లేదా
$x=1$ అయితే $2\leq\fact{x}+1$, కనుక వాదన తక్షణమే సిద్ధిస్తుంది.
ఇక $x\geq2$ అనుకుందాం. $p_n$ అనేది $\le x$ అయిన అతి పెద్ద ప్రధాన
సంఖ్య అనుకొని, $\le x$ అయిన ప్రధాన సంఖ్యలన్నిటి లబ్ధం
$p=p_0\cdot p_1\cdot\dots\cdot p_n$ను పరిశీలించండి. $p+1$ ప్రధాన
సంఖ్య అయి ఉండాలి, లేదా $x$, $p+1$ మధ్య ఒక ప్రధాన సంఖ్య ఉండాలి. ఎందుకు?
$p+1$ ప్రధాన సంఖ్య కాదనుకుందాం. అప్పుడు $q<p+1$ అయ్యే ఏదో ప్రధాన
సంఖ్య $q\mid p+1$. $\le x$ అయిన ప్రధాన సంఖ్యల్లో ఏదీ $p+1$ను
భాగించదు. ($p$ నిర్వచనం వల్ల ప్రతి $p_i\le x$, $p$ను శేషం~$0$తో
భాగిస్తుంది. అందువల్ల ప్రతి $p_i\le x$, $p+1$ను శేషం~$1$తో
భాగిస్తుంది; కాబట్టి $p_i\nmid p+1$.) కనుక $q$ అనే ప్రధాన సంఖ్య
$>x$, $<p+1$. అంతేకాక $p\le\fact{x}$; అందువల్ల $>x$, $\le
\fact{x}+1$ అయిన ప్రధాన సంఖ్య ఒకటి ఉంది.)
\sourcecorrection{OLTECRREM-005}{ఆంగ్ల నిరూపణ $p_n$ను ``$\leq x$ అయిన అతి పెద్ద ప్రధాన సంఖ్య''గా తీసుకున్నప్పటికీ $x=0,1$ సందర్భాలను విడిగా చూడలేదు. ఆ రెండు సందర్భాల్లో ప్రధాన సంఖ్య $2$ నేరుగా అవసరమైన అవధిలో ఉంటుందని చేర్చి, లబ్ధ వాదనను $x\geq2$కు పరిమితం చేశాం.}

\begin{prob}
పరిమిత కనిష్ఠీకరణను ఉపయోగించి పూర్ణాంక భాగహారం $d(x,y)$ను
నిర్వచించండి.
\end{prob}

\end{document}
